% !TeX root = ../main.tex

\begin{abstract}

本論文針對上行非正交的多重存取（non-orthogonal multiple access, NOMA）系統進行比較分析。首先，在未編碼單一資源系統中，比較功率域非正交的多重存取（power-domain NOMA, PD-NOMA）與分增益多重存取（gain division multiple access, GDMA）的錯誤率表現。
接著將功率域非正交的多重存取 (PD-NOMA) 延伸至多資源系統模型，並在未編碼傳輸下比較 MMSE-SIC 偵測器與高密度簽名碼分
碼多重存取（high density signature CDMA, HDS-CDMA）的效能。

其次，本論文探討不同的編碼多資源 NOMA 接收器，並進一步將這些接收器延伸至正交分頻多工（orthogonal frequency division multiplexing, OFDM）架構下的 NOMA 系統。

最後，我們提出一個基於調變編碼方式（modulation coding scheme, MCS）的使用者配對架構，用以評估 NOMA 資源共享對系統層級頻譜效率（spectral efficiency, SE）表現的影響。

\end{abstract}

\begin{abstract*}

This thesis presents a comparative study of uplink non-orthogonal multiple access (NOMA) systems. 
First, uncoded single-resource power-domain NOMA (PD-NOMA) is compared with gain division multiple access (GDMA). PD-NOMA is then extended to a multi-resource model, where MMSE-SIC detection is evaluated against high density signature CDMA (HDS-CDMA) under uncoded transmission. 

Second, different coded multi-resource NOMA receivers are discussed. This receivers are further extended to orthogonal frequency division multiplexing (OFDM)-based NOMA systems.

Finally, a modulation coding scheme (MCS)-based user-pairing framework is developed to evaluate the system-level spectral efficiency (SE) behavior caused by NOMA resource sharing.

\end{abstract*}